\documentclass[11pt,a4paper]{article}

\usepackage[a4paper, top=0.9in, bottom=0.9in, left=1in, right=1in]{geometry}
\usepackage[utf8]{inputenc}
\usepackage[T1]{fontenc}
\usepackage{charter}
\usepackage{parskip}

\setlength{\parskip}{0.8em}
\setlength{\parindent}{0pt}
\pagestyle{empty}

\begin{document}

Dear Hiring Team,

I am applying for the Python Developer role at SocietyWorks. I use
TheyWorkForYou and WhatDoTheyKnow regularly, and earlier this month I
reached out to volunteer on the register-of-interests scrutiny project. I
am also a software engineer who has independently built an open-source
project that makes UK wealth inequality data accessible to the public, a
choice that grew from my widening participation work running workshops with
secondary school students who talk about feeling the weight of economic
pressures they cannot see. Making invisible systems visible through data
and open-source tools is what mySociety does at scale. It is also what I
have been building on my own.

My strongest Python work is on that project: a FastAPI backend deployed via
GitHub Actions with 8 CI/CD pipelines, 11 reproducible data pipelines
ingesting ONS, HMRC, and WID government datasets, a policy
microsimulation prototype for modelling potential UK tax reforms, and 72
Vue~3 components including 13 interactive chart visualisations with
WCAG~AA accessibility. I also serve as DevOps and Security Lead of a
FastAPI platform built with an international team of four, where I
maintain 120+ API endpoints and 39{,}500+ automated tests, and led a
security hardening sprint remediating SSRF, SQL injection, and other
CRITICAL vulnerabilities across 640+ reviewed pull requests.

During my 15-month internship at GE Digital, I integrated automated Nmap
scanning into approximately 72 AWS CI/CD pipelines using Jenkins, Bash,
and Groovy, and received a GE Spotlight Award for the work. For the Perl
aspect of this role: I also maintained a legacy Smallworld/Magik codebase,
a niche language I had never seen before joining. I learned it on the job,
fixed bugs, improved database interactions, added unit tests, and wrote
onboarding documentation. I am confident I can do the same with Perl. I
work comfortably with Docker, GitHub Actions, and Linux, and have hands-on
experience with infrastructure around busy web services from both the
enterprise environment at GE and my own deployment work.

In my university role, I deliver widening participation talks and
workshops, regularly explaining technical ideas to non-technical audiences.
This translates directly to client-facing support, demos, and cross-team
collaboration. I have also taken on a data role well beyond my job
description, building Python scripts for data cleaning, visualisation, and
process automation for the team, and I am now developing a programme to
bring data literacy workshops into schools. I have presented as a
conference speaker at SGAI-AI 2025, where I published a lead-author paper
in Springer proceedings. I am disciplined about code review and
collaborative development, and I write code that others can maintain,
because my projects are designed for open-source contributors.

I would be genuinely excited to bring my Python, web, and infrastructure
skills to services used by over 30 million people, to learn Perl alongside
experienced colleagues, and to grow within a team whose civic mission I
already share. I am comfortable joining an on-call rota and
context-switching between development, maintenance, and support.

\vspace{0.3em}
C.T.

\end{document}
